% ===========================================================================
%  evnx for Full-stack / Java Developers  —  Spring Boot + Vue + React monorepo
%  Persona RRP.  Build: pdflatex -interaction=nonstopmode fullstack-java.tex
% ===========================================================================
\documentclass[aspectratio=169,10pt]{beamer}

\newcommand{\capturedir}{../captures}
% ---------------------------------------------------------------------------
%  Shared preamble for the seven audience decks.
%
%  Each deck does, before the document body starts:
%     \newcommand{\capturedir}{../captures}
%     % ---------------------------------------------------------------------------
%  Shared preamble for the seven audience decks.
%
%  Each deck does, before the document body starts:
%     \newcommand{\capturedir}{../captures}
%     % ---------------------------------------------------------------------------
%  Shared preamble for the seven audience decks.
%
%  Each deck does, before the document body starts:
%     \newcommand{\capturedir}{../captures}
%     \input{persona-preamble}
%
%  \capturedir points at ../captures so both pools resolve:
%     \capture{24-scan.txt}              — the original v0.5.0 captures
%     \capture{persona/fe-01-scan-dist.txt} — captures recorded for these decks
% ---------------------------------------------------------------------------

\input{../common/evnx-preamble}

% ---- provenance ------------------------------------------------------------
% The original decks were pinned to tag v0.5.0 (85a52c3). Everything captured
% for the persona decks was re-run against 0.5.2, so the footline differs.
\renewcommand{\verifiedon}{\tiny\textcolor{evnxMuted}{verified against evnx 0.5.2 --- real run, not transcribed}}

% ---- audience badge on the title slide ------------------------------------
\newcommand{\audience}[1]{%
  \begin{center}\small\textcolor{evnxAccent}{\textbf{Audience: #1}}\end{center}}

% ---- verdict chips ---------------------------------------------------------
\newcommand{\solved}{\textcolor{evnxOk}{\ensuremath{\checkmark}}}
\newcommand{\partialmark}{\textcolor{evnxWarn}{\textbf{\textasciitilde}}}
\newcommand{\unsolved}{\textcolor{evnxErr}{\ensuremath{\times}}}

% ---- a correction block: something the source material got wrong -----------
\newenvironment{correction}[1]{%
  \begin{alertblock}{\small Correction --- #1}\small}{\end{alertblock}}

% ---- a trap block: works in the demo, fails in real use --------------------
\newenvironment{trap}[1]{%
  \setbeamercolor{block title}{bg=evnxWarn,fg=white}%
  \setbeamercolor{block body}{bg=evnxWarn!8}%
  \begin{block}{\small #1}\small}{\end{block}}

% ---- the persona card ------------------------------------------------------
\newcommand{\personacard}[4]{%
  \begin{block}{#1}
    \small
    \textbf{Stack:} #2 \\
    \textbf{Ships:} #3 \\
    \textbf{Loses time to:} #4
  \end{block}}

% ---- speaker note (rendered small and muted, not hidden) -------------------
\newcommand{\saythis}[1]{%
  \vspace{0.4em}{\footnotesize\textcolor{evnxMuted}{\textit{Say:} #1}}}

%
%  \capturedir points at ../captures so both pools resolve:
%     \capture{24-scan.txt}              — the original v0.5.0 captures
%     \capture{persona/fe-01-scan-dist.txt} — captures recorded for these decks
% ---------------------------------------------------------------------------

% ---------------------------------------------------------------------------
%  evnx presentation preamble — shared by all three decks
%
%  Engine: pdfLaTeX (also builds under XeLaTeX / LuaLaTeX).
%  Packages: beamer, listings, xcolor, amssymb, newunicodechar, booktabs,
%            fancyvrb, hyperref — all in texlive-latex-recommended + beamer.
%
%  \capturedir must be defined by the including document BEFORE \input of
%  this file, e.g.  \newcommand{\capturedir}{../captures}
% ---------------------------------------------------------------------------

\usetheme{Madrid}
\usecolortheme{seahorse}
\usefonttheme{professionalfonts}
\setbeamertemplate{navigation symbols}{}
\setbeamertemplate{caption}[numbered]
\setbeamerfont{frametitle}{size=\large}
\setbeamerfont{title}{size=\LARGE,series=\bfseries}

\usepackage[T1]{fontenc}
\usepackage[utf8]{inputenc}
\usepackage{lmodern}
\usepackage{amssymb}
\usepackage{booktabs}
\usepackage{xcolor}
\usepackage{listings}
\usepackage{fancyvrb}
\usepackage{newunicodechar}
\usepackage{tikz}
\usetikzlibrary{arrows.meta,positioning,shapes.geometric,fit,backgrounds}

% ---- palette ---------------------------------------------------------------
\definecolor{evnxInk}{HTML}{1F2933}
\definecolor{evnxAccent}{HTML}{2F6F4E}
\definecolor{evnxWarn}{HTML}{B7791F}
\definecolor{evnxErr}{HTML}{B03A2E}
\definecolor{evnxOk}{HTML}{2F6F4E}
\definecolor{evnxMuted}{HTML}{6B7280}
\definecolor{evnxTermBg}{HTML}{F5F5F2}
\definecolor{evnxRule}{HTML}{D1D5DB}

\setbeamercolor{structure}{fg=evnxAccent}
\setbeamercolor{alerted text}{fg=evnxErr}

% ---- glyph mapping for body text ------------------------------------------
% evnx emits exactly these 11 non-ASCII characters (enumerated from the
% captured logs, not guessed).
\newcommand{\evnxCheck}{\textcolor{evnxOk}{\ensuremath{\checkmark}}}
\newcommand{\evnxCross}{\textcolor{evnxErr}{\ensuremath{\times}}}
\newcommand{\evnxWarnSign}{\textcolor{evnxWarn}{\textbf{!}}}

\newunicodechar{✓}{\evnxCheck}
\newunicodechar{✗}{\evnxCross}
\newunicodechar{⚠}{\evnxWarnSign}
\newunicodechar{→}{\ensuremath{\rightarrow}}
\newunicodechar{↔}{\ensuremath{\leftrightarrow}}
\newunicodechar{─}{\textendash}
\newunicodechar{—}{---}
\newunicodechar{·}{\textperiodcentered}
\newunicodechar{…}{\ldots}
\newunicodechar{•}{\textbullet}
\newunicodechar{️}{}                % U+FE0F variation selector — invisible

% ---- terminal listing style ------------------------------------------------
% `literate` is required: newunicodechar does not apply inside listings.
\lstdefinestyle{evnxterm}{
  basicstyle=\ttfamily\scriptsize,
  backgroundcolor=\color{evnxTermBg},
  frame=single,
  rulecolor=\color{evnxRule},
  framesep=4pt,
  breaklines=true,
  breakatwhitespace=false,
  columns=fullflexible,
  keepspaces=true,
  showstringspaces=false,
  upquote=true,
  extendedchars=true,
  literate=
    {✓}{{\textcolor{evnxOk}{$\checkmark$}}}1
    {✗}{{\textcolor{evnxErr}{$\times$}}}1
    {⚠}{{\textcolor{evnxWarn}{\textbf{!}}}}1
    {→}{{$\rightarrow$}}1
    {↔}{{$\leftrightarrow$}}1
    {─}{{-}}1
    {—}{{---}}1
    {·}{{$\cdot$}}1
    {…}{{\ldots}}1
    {•}{{\textbullet}}1
    {️}{{}}0
    {é}{{\'e}}1
}

\lstdefinestyle{evnxcode}{
  style=evnxterm,
  basicstyle=\ttfamily\scriptsize,
  backgroundcolor=\color{white},
}

\lstset{style=evnxterm}

% ---- helper macros ---------------------------------------------------------
% \capture{file}        — include a real captured run verbatim
% \capturepart{f}{a}{b} — include lines a..b of a captured run
\newcommand{\capture}[1]{\lstinputlisting[style=evnxterm]{\capturedir/#1}}
\newcommand{\capturepart}[3]{%
  \lstinputlisting[style=evnxterm,firstline=#2,lastline=#3]{\capturedir/#1}}

% a labelled "real output" box
\newcommand{\realrun}[1]{%
  \begin{center}\scriptsize\textcolor{evnxMuted}{real run \textendash\ #1}\end{center}}

\newcommand{\cmd}[1]{\texttt{\textbf{#1}}}
\newcommand{\flag}[1]{\texttt{#1}}
\newcommand{\file}[1]{\texttt{#1}}

% exit-code chips
\newcommand{\exitzero}{\textcolor{evnxOk}{\texttt{0}}}
\newcommand{\exitone}{\textcolor{evnxWarn}{\texttt{1}}}
\newcommand{\exittwo}{\textcolor{evnxErr}{\texttt{2}}}

% a standard "verified against" footline for factual slides
\newcommand{\verifiedon}{\tiny\textcolor{evnxMuted}{verified against evnx 0.5.0 @ 9cfe75b}}

\usepackage[colorlinks=true,linkcolor=evnxAccent,urlcolor=evnxAccent]{hyperref}


% ---- provenance ------------------------------------------------------------
% The original decks were pinned to tag v0.5.0 (85a52c3). Everything captured
% for the persona decks was re-run against 0.5.2, so the footline differs.
\renewcommand{\verifiedon}{\tiny\textcolor{evnxMuted}{verified against evnx 0.5.2 --- real run, not transcribed}}

% ---- audience badge on the title slide ------------------------------------
\newcommand{\audience}[1]{%
  \begin{center}\small\textcolor{evnxAccent}{\textbf{Audience: #1}}\end{center}}

% ---- verdict chips ---------------------------------------------------------
\newcommand{\solved}{\textcolor{evnxOk}{\ensuremath{\checkmark}}}
\newcommand{\partialmark}{\textcolor{evnxWarn}{\textbf{\textasciitilde}}}
\newcommand{\unsolved}{\textcolor{evnxErr}{\ensuremath{\times}}}

% ---- a correction block: something the source material got wrong -----------
\newenvironment{correction}[1]{%
  \begin{alertblock}{\small Correction --- #1}\small}{\end{alertblock}}

% ---- a trap block: works in the demo, fails in real use --------------------
\newenvironment{trap}[1]{%
  \setbeamercolor{block title}{bg=evnxWarn,fg=white}%
  \setbeamercolor{block body}{bg=evnxWarn!8}%
  \begin{block}{\small #1}\small}{\end{block}}

% ---- the persona card ------------------------------------------------------
\newcommand{\personacard}[4]{%
  \begin{block}{#1}
    \small
    \textbf{Stack:} #2 \\
    \textbf{Ships:} #3 \\
    \textbf{Loses time to:} #4
  \end{block}}

% ---- speaker note (rendered small and muted, not hidden) -------------------
\newcommand{\saythis}[1]{%
  \vspace{0.4em}{\footnotesize\textcolor{evnxMuted}{\textit{Say:} #1}}}

%
%  \capturedir points at ../captures so both pools resolve:
%     \capture{24-scan.txt}              — the original v0.5.0 captures
%     \capture{persona/fe-01-scan-dist.txt} — captures recorded for these decks
% ---------------------------------------------------------------------------

% ---------------------------------------------------------------------------
%  evnx presentation preamble — shared by all three decks
%
%  Engine: pdfLaTeX (also builds under XeLaTeX / LuaLaTeX).
%  Packages: beamer, listings, xcolor, amssymb, newunicodechar, booktabs,
%            fancyvrb, hyperref — all in texlive-latex-recommended + beamer.
%
%  \capturedir must be defined by the including document BEFORE \input of
%  this file, e.g.  \newcommand{\capturedir}{../captures}
% ---------------------------------------------------------------------------

\usetheme{Madrid}
\usecolortheme{seahorse}
\usefonttheme{professionalfonts}
\setbeamertemplate{navigation symbols}{}
\setbeamertemplate{caption}[numbered]
\setbeamerfont{frametitle}{size=\large}
\setbeamerfont{title}{size=\LARGE,series=\bfseries}

\usepackage[T1]{fontenc}
\usepackage[utf8]{inputenc}
\usepackage{lmodern}
\usepackage{amssymb}
\usepackage{booktabs}
\usepackage{xcolor}
\usepackage{listings}
\usepackage{fancyvrb}
\usepackage{newunicodechar}
\usepackage{tikz}
\usetikzlibrary{arrows.meta,positioning,shapes.geometric,fit,backgrounds}

% ---- palette ---------------------------------------------------------------
\definecolor{evnxInk}{HTML}{1F2933}
\definecolor{evnxAccent}{HTML}{2F6F4E}
\definecolor{evnxWarn}{HTML}{B7791F}
\definecolor{evnxErr}{HTML}{B03A2E}
\definecolor{evnxOk}{HTML}{2F6F4E}
\definecolor{evnxMuted}{HTML}{6B7280}
\definecolor{evnxTermBg}{HTML}{F5F5F2}
\definecolor{evnxRule}{HTML}{D1D5DB}

\setbeamercolor{structure}{fg=evnxAccent}
\setbeamercolor{alerted text}{fg=evnxErr}

% ---- glyph mapping for body text ------------------------------------------
% evnx emits exactly these 11 non-ASCII characters (enumerated from the
% captured logs, not guessed).
\newcommand{\evnxCheck}{\textcolor{evnxOk}{\ensuremath{\checkmark}}}
\newcommand{\evnxCross}{\textcolor{evnxErr}{\ensuremath{\times}}}
\newcommand{\evnxWarnSign}{\textcolor{evnxWarn}{\textbf{!}}}

\newunicodechar{✓}{\evnxCheck}
\newunicodechar{✗}{\evnxCross}
\newunicodechar{⚠}{\evnxWarnSign}
\newunicodechar{→}{\ensuremath{\rightarrow}}
\newunicodechar{↔}{\ensuremath{\leftrightarrow}}
\newunicodechar{─}{\textendash}
\newunicodechar{—}{---}
\newunicodechar{·}{\textperiodcentered}
\newunicodechar{…}{\ldots}
\newunicodechar{•}{\textbullet}
\newunicodechar{️}{}                % U+FE0F variation selector — invisible

% ---- terminal listing style ------------------------------------------------
% `literate` is required: newunicodechar does not apply inside listings.
\lstdefinestyle{evnxterm}{
  basicstyle=\ttfamily\scriptsize,
  backgroundcolor=\color{evnxTermBg},
  frame=single,
  rulecolor=\color{evnxRule},
  framesep=4pt,
  breaklines=true,
  breakatwhitespace=false,
  columns=fullflexible,
  keepspaces=true,
  showstringspaces=false,
  upquote=true,
  extendedchars=true,
  literate=
    {✓}{{\textcolor{evnxOk}{$\checkmark$}}}1
    {✗}{{\textcolor{evnxErr}{$\times$}}}1
    {⚠}{{\textcolor{evnxWarn}{\textbf{!}}}}1
    {→}{{$\rightarrow$}}1
    {↔}{{$\leftrightarrow$}}1
    {─}{{-}}1
    {—}{{---}}1
    {·}{{$\cdot$}}1
    {…}{{\ldots}}1
    {•}{{\textbullet}}1
    {️}{{}}0
    {é}{{\'e}}1
}

\lstdefinestyle{evnxcode}{
  style=evnxterm,
  basicstyle=\ttfamily\scriptsize,
  backgroundcolor=\color{white},
}

\lstset{style=evnxterm}

% ---- helper macros ---------------------------------------------------------
% \capture{file}        — include a real captured run verbatim
% \capturepart{f}{a}{b} — include lines a..b of a captured run
\newcommand{\capture}[1]{\lstinputlisting[style=evnxterm]{\capturedir/#1}}
\newcommand{\capturepart}[3]{%
  \lstinputlisting[style=evnxterm,firstline=#2,lastline=#3]{\capturedir/#1}}

% a labelled "real output" box
\newcommand{\realrun}[1]{%
  \begin{center}\scriptsize\textcolor{evnxMuted}{real run \textendash\ #1}\end{center}}

\newcommand{\cmd}[1]{\texttt{\textbf{#1}}}
\newcommand{\flag}[1]{\texttt{#1}}
\newcommand{\file}[1]{\texttt{#1}}

% exit-code chips
\newcommand{\exitzero}{\textcolor{evnxOk}{\texttt{0}}}
\newcommand{\exitone}{\textcolor{evnxWarn}{\texttt{1}}}
\newcommand{\exittwo}{\textcolor{evnxErr}{\texttt{2}}}

% a standard "verified against" footline for factual slides
\newcommand{\verifiedon}{\tiny\textcolor{evnxMuted}{verified against evnx 0.5.0 @ 9cfe75b}}

\usepackage[colorlinks=true,linkcolor=evnxAccent,urlcolor=evnxAccent]{hyperref}


% ---- provenance ------------------------------------------------------------
% The original decks were pinned to tag v0.5.0 (85a52c3). Everything captured
% for the persona decks was re-run against 0.5.2, so the footline differs.
\renewcommand{\verifiedon}{\tiny\textcolor{evnxMuted}{verified against evnx 0.5.2 --- real run, not transcribed}}

% ---- audience badge on the title slide ------------------------------------
\newcommand{\audience}[1]{%
  \begin{center}\small\textcolor{evnxAccent}{\textbf{Audience: #1}}\end{center}}

% ---- verdict chips ---------------------------------------------------------
\newcommand{\solved}{\textcolor{evnxOk}{\ensuremath{\checkmark}}}
\newcommand{\partialmark}{\textcolor{evnxWarn}{\textbf{\textasciitilde}}}
\newcommand{\unsolved}{\textcolor{evnxErr}{\ensuremath{\times}}}

% ---- a correction block: something the source material got wrong -----------
\newenvironment{correction}[1]{%
  \begin{alertblock}{\small Correction --- #1}\small}{\end{alertblock}}

% ---- a trap block: works in the demo, fails in real use --------------------
\newenvironment{trap}[1]{%
  \setbeamercolor{block title}{bg=evnxWarn,fg=white}%
  \setbeamercolor{block body}{bg=evnxWarn!8}%
  \begin{block}{\small #1}\small}{\end{block}}

% ---- the persona card ------------------------------------------------------
\newcommand{\personacard}[4]{%
  \begin{block}{#1}
    \small
    \textbf{Stack:} #2 \\
    \textbf{Ships:} #3 \\
    \textbf{Loses time to:} #4
  \end{block}}

% ---- speaker note (rendered small and muted, not hidden) -------------------
\newcommand{\saythis}[1]{%
  \vspace{0.4em}{\footnotesize\textcolor{evnxMuted}{\textit{Say:} #1}}}


\title{evnx for full-stack teams}
\subtitle{Spring Boot, Vue and React in one repo --- three prefixes, one URL}
\date{}

\begin{document}

\begin{frame}[plain]
  \titlepage
  \audience{Spring Boot 3 · Vue 3 + Vite · React (CRA) · monorepo · Kubernetes}
  \begin{center}\footnotesize
  \textcolor{evnxMuted}{evnx 0.5.2. Every terminal block is a real run.}
  \end{center}
\end{frame}

% =========================================================================
\begin{frame}{This deck in one slide}
  \begin{block}{Your daily annoyance has a clean fix}
    Config trapped in someone's IntelliJ run configuration, onboarded by
    screenshot. \cmd{evnx cloud run -{}- ./mvnw spring-boot:run} ends that,
    today, with no change to your application code.
  \end{block}

  \vspace{0.7em}
  \begin{block}{Your structural problem does not}
    Three apps, three prefixes, one API URL defined three times. evnx has
    \textbf{no monorepo model and no cross-vault references}. You will run it
    per directory and keep the duplication --- or generate the prefixed names
    from one base, which we will demo.
  \end{block}
\end{frame}

\begin{frame}{Contents}\tableofcontents[hideallsubsections]\end{frame}

% =========================================================================
\section{A day in the monorepo}

\begin{frame}{Three things, one afternoon}
  \small
  \begin{block}{The new laptop}
    The backend will not start: \texttt{DB\_PASSWORD} lives in a colleague's
    IntelliJ run configuration, which is not in git. You copy twelve variables
    by hand from a screenshot.
  \end{block}

  \begin{block}{The URL that moved}
    You update \file{backend/application-dev.yml},
    \file{admin-ui/.env.development} (\texttt{VITE\_API\_URL}) and
    \file{customer-web/.env.development} (\texttt{REACT\_APP\_API\_URL}). You
    miss the React one. QA files a bug.
  \end{block}

  \begin{block}{The year-old commit}
    Security scanning flags
    \texttt{spring.datasource.password: S3cr3t!} committed in
    \file{application-staging.yml}.
  \end{block}
\end{frame}

\begin{frame}{Scored}
  \small
  \begin{tabular}{@{}llp{6.2cm}@{}}
    \toprule
    \textbf{Problem} & & \textbf{evnx} \\
    \midrule
    IntelliJ / screenshot onboarding & \solved & vault + \cmd{cloud run}. The clean win \\
    Spring does not read \file{.env}  & \solved & \cmd{cloud run} injects env; relaxed binding does the rest \\
    One URL, three names             & \partialmark & \flag{-{}-prefix} generates the names from one base \\
    Secrets in \file{application-*.yml} & \solved & \cmd{scan} reads YAML, not just \file{.env} \\
    Keystores, PEM, multi-line JSON  & \partialmark & \textbf{revised} --- see slide 14 \\
    \bottomrule
  \end{tabular}

  \vspace{0.9em}
  \saythis{``Four of five have an answer today. Let me show you each one
  running, and then the one that does not.''}
\end{frame}

% =========================================================================
\section{Spring without .env support}

\begin{frame}[fragile]{\cmd{cloud run} --- and relaxed binding does the mapping}
\begin{lstlisting}[style=evnxcode]
cd backend && evnx cloud run -- ./mvnw spring-boot:run
\end{lstlisting}

  \vspace{0.6em}
  Spring Boot's relaxed binding maps environment variables onto properties
  automatically:

  \vspace{0.4em}
  \begin{center}\small
  \begin{tabular}{@{}ll@{}}
    \toprule
    \textbf{Vault variable} & \textbf{Spring property} \\
    \midrule
    \texttt{SPRING\_DATASOURCE\_URL}   & \texttt{spring.datasource.url} \\
    \texttt{SPRING\_DATASOURCE\_PASSWORD} & \texttt{spring.datasource.password} \\
    \texttt{APP\_PAYMENT\_API\_KEY}    & \texttt{app.payment.api-key} \\
    \bottomrule
  \end{tabular}
  \end{center}

  \vspace{0.6em}
  {\footnotesize No \texttt{spring.config.import}, no EnvFile plugin, no code
  change. If you prefer files, \texttt{spring.config.import=optional:file:.env[.properties]}
  also works --- but then you are back to a plaintext file on disk.}
\end{frame}

\begin{frame}[fragile]{The YAML stays clean}
\begin{lstlisting}[style=evnxcode]
spring:
  datasource:
    url: ${DB_URL}
    password: ${DB_PASSWORD}
app:
  payment:
    api-key: ${APP_PAYMENT_API_KEY}
\end{lstlisting}

  \vspace{0.7em}
  \begin{block}{Same file, three sources}
    Locally the values come from \cmd{cloud run}. In CI from the runner's
    environment. In production from Kubernetes. The committed YAML never
    changes and never holds a secret.
  \end{block}

  \vspace{0.5em}
  {\footnotesize Also covers Testcontainers and Kafka integration tests:
  \cmd{evnx cloud run -{}- ./mvnw verify}.}
\end{frame}

\begin{frame}[fragile]{\cmd{scan} reads your resources directory}
  \capturepart{persona/fs-01-scan-resources.txt}{1}{14}
  \realrun{\texttt{evnx scan backend/src/main/resources/}}

  \vspace{0.4em}
  \begin{correction}{\cmd{scan} is not a \file{.env}-only tool}
    It reads \texttt{yml}, \texttt{yaml}, \texttt{java}, \texttt{properties},
    \texttt{xml}, \texttt{json}, \texttt{ts}, \texttt{js}, \texttt{py},
    \texttt{rs}, \texttt{go} and more. A widespread assumption says it only
    looks at \file{.env} --- it does not, and that is what makes it useful on a
    Spring codebase.
  \end{correction}
\end{frame}

% =========================================================================
\section{One value, three names}

\begin{frame}[fragile]{\flag{-{}-prefix} generates the per-app names}
  \begin{columns}[T]
    \begin{column}{0.49\textwidth}
      \capture{persona/fs-02-prefix-vite.txt}
    \end{column}
    \begin{column}{0.49\textwidth}
      \capture{persona/fs-03-prefix-cra.txt}
    \end{column}
  \end{columns}
  \realrun{one \file{.env}, two prefixed outputs}

  \vspace{0.4em}
  {\footnotesize Keep \textbf{one} base value. Generate
  \texttt{VITE\_API\_URL} and \texttt{REACT\_APP\_API\_URL} from it in each
  app's build step. The QA bug from slide 5 becomes impossible because there is
  only one place to edit.}
\end{frame}

\begin{frame}[fragile]{Wiring it into the monorepo}
\begin{lstlisting}[style=evnxcode]
# admin-ui/prebuild.sh
evnx convert --env ../shared.env --to shell --prefix VITE_ > .env.generated

# customer-web/prebuild.sh
evnx convert --env ../shared.env --to shell --prefix REACT_APP_ > .env.generated
\end{lstlisting}

  \vspace{0.7em}
  \begin{trap}{This is a pattern you maintain, not a feature}
    There is no \texttt{evnx.toml} listing workspace directories, no
    \cmd{validate -{}-all}, no \cmd{cloud pull -{}-all}, and no
    \texttt{\$\{\{shared.API\_URL\}\}} reference syntax across vaults. Monorepo
    workspaces are item 1 on this audience's roadmap.
  \end{trap}

  \vspace{0.4em}
  {\footnotesize Until then: three vaults (\texttt{shop-backend},
  \texttt{shop-admin}, \texttt{shop-web}) and a CI matrix over the three
  directories.}
\end{frame}

\begin{frame}[fragile]{Also useful: the CRA → Vite migration}
\begin{lstlisting}[style=evnxcode]
# what needs renaming?
evnx convert --to json --include "REACT_APP_*" | jq -r 'keys[]'
\end{lstlisting}

  \vspace{0.7em}
  \begin{block}{No rename command}
    \cmd{convert} lists and re-prefixes; it does not rewrite your source. Pair
    the list above with \texttt{sed} over the codebase, then
    \cmd{evnx diff} to confirm nothing was dropped.
  \end{block}
\end{frame}

% =========================================================================
\section{The frontend halves}

\begin{frame}[fragile]{\texttt{VITE\_} and \texttt{REACT\_APP\_} leak the same way}
  \capture{persona/fs-04-scan-cra-build.txt}
  \realrun{\file{customer-web/build/static/js/main.js} contains a live Stripe key}

  \vspace{0.4em}
  \begin{trap}{\file{build/} and \file{dist/} are on the default exclusion list}
    The standard advice --- \cmd{evnx scan customer-web/build/} --- reports
    ``0 files scanned'' and exit \exitzero. \file{src/commands/scan/filters.rs:106}.
    Tracked \textbf{NEW-1, P0}.
  \end{trap}

  \vspace{0.4em}
  {\footnotesize Workaround: \texttt{cp -r build /tmp/scanme \&\& evnx scan
  /tmp/scanme}. And assert \texttt{.files\_scanned > 0} in CI rather than
  trusting the exit code.}
\end{frame}

\begin{frame}{The build-time trap, stated once for both frameworks}
  \small
  \begin{itemize}
    \item \texttt{VITE\_*} and \texttt{REACT\_APP\_*} are \textbf{inlined at
          build time}. Changing the vault does nothing until you rebuild.
    \item Frontend Docker images therefore need the values as \textbf{build
          args} --- \cmd{convert -{}-to shell} feeds them --- or you move to a
          runtime \file{env.js} rendered by \cmd{evnx template} at container
          start.
    \item evnx has \textbf{no prefix awareness}: a secret named
          \texttt{VITE\_STRIPE\_KEY} passes \cmd{validate} cleanly. Use
          \cmd{scan}, which works on values, plus a name guard of your own.
  \end{itemize}

  \vspace{0.7em}
  \begin{block}{The name guard, one line}
    \begin{center}\ttfamily\scriptsize
    evnx convert -{}-to json \textbar{} jq -r 'keys[] \textbar{} select(test("\^{}(VITE\_\textbar{}REACT\_APP\_)")) \textbar{} select(test("SECRET\textbar{}TOKEN\textbar{}KEY"))'
    \end{center}
  \end{block}
\end{frame}

% =========================================================================
\section{Multi-line and binary}

\begin{frame}{The claim to correct}
  \begin{center}\Large
  ``Keystores, PEM certificates, multi-line JSON --- \\
  \textbf{parser does not support multi-line values}.''
  \end{center}

  \vspace{1.0em}
  \begin{center}\large
  Half wrong. Multi-line values parse. \emph{Binary files} are the real gap.
  \end{center}

  \vspace{0.8em}
  \saythis{Splitting the claim is the point --- one half is a solved problem
  people are working around unnecessarily.}
\end{frame}

\begin{frame}[fragile]{A PEM certificate, round-tripped}
  \capture{persona/py-05-multiline-pem.txt}
  \realrun{double-quoted multi-line value}

  \vspace{0.4em}
  {\footnotesize \texttt{B=2} on the line after proves the parser resumed
  correctly. Newlines survive as \texttt{\textbackslash n}. Spring reads it
  from the environment exactly as written.}
\end{frame}

\begin{frame}[fragile]{The quoting rule that decides whether it works}
  \begin{columns}[T]
    \begin{column}{0.49\textwidth}
      \centering\footnotesize\textbf{single-quoted JSON --- works}
      \capturepart{persona/py-07-multiline-json-sq.txt}{1}{8}
    \end{column}
    \begin{column}{0.49\textwidth}
      \centering\footnotesize\textbf{double-quoted with \texttt{\textbackslash "} --- fails}
      \capturepart{persona/py-06-multiline-json-dq.txt}{1}{5}
    \end{column}
  \end{columns}

  \vspace{0.6em}
  \begin{block}{The rule}
    A double-quoted value ends at the first inner \texttt{"}, escaped or not.
    \textbf{Wrap JSON in single quotes.} PEM has no inner quotes so either
    works. The failure is loud --- exit \exittwo\ with a line number --- which
    is the right kind of failure.
  \end{block}
\end{frame}

\begin{frame}{What genuinely does not fit in a vault}
  \small
  \begin{tabular}{@{}lll@{}}
    \toprule
    \textbf{Artefact} & & \textbf{Today} \\
    \midrule
    PEM certificate, inline          & \solved   & a multi-line variable \\
    Service-account JSON, inline     & \solved   & single-quoted multi-line \\
    Java keystore \file{.jks}, \file{.p12} & \unsolved & binary. Mount it; vault the password \\
    Spring Cloud Config / HashiCorp Vault & \unsolved & no sync in either direction \\
    Actuator \texttt{/env} exposure  & \unsolved & Spring configuration, not a file problem \\
    \bottomrule
  \end{tabular}

  \vspace{0.8em}
  \begin{block}{If your company already runs HashiCorp Vault}
    Decide the source of truth \textbf{per environment}, not per tool. The
    workable split is evnx for local and dev, Vault for staging and production,
    with \cmd{convert} / \cmd{migrate} as the hand-off. Two systems pretending
    to both own production is the failure mode.
  \end{block}
\end{frame}

% =========================================================================
\section{Deploy and CI}

\begin{frame}[fragile]{Kubernetes, and the matrix build}
\begin{lstlisting}[style=evnxcode]
# K8s Secret -- pipe it, never commit it
evnx convert --to kubernetes | kubectl apply -f -

# CI: one job per app directory
strategy:
  matrix:
    app: [backend, admin-ui, customer-web]
steps:
  - run: evnx validate --strict --format github
    working-directory: ${{ matrix.app }}
  - run: evnx scan . --format sarif > evnx.sarif || test $? -eq 1
\end{lstlisting}

  \vspace{0.5em}
  {\footnotesize Helm values: \cmd{evnx template} renders them. There is no
  Helm-native output target and no Argo CD integration.}
\end{frame}

\begin{frame}{The scorecard, revised}
  \begin{center}
  \begin{tabular}{@{}clp{6.2cm}@{}}
    \toprule
    & \textbf{Count} & \textbf{Items} \\
    \midrule
    \solved     & 7 & IntelliJ onboarding, Spring via \cmd{cloud run}, \file{.env.example} drift, clean YAML, secret scanning in resources, CI gates, \textbf{multi-line values} \\
    \partialmark & 9 & monorepo structure, one URL three names, relaxed-binding naming, frontend prefix leaks, integration tests, CRA migration, Vault overlap, rotation, K8s/Helm \\
    \unsolved   & 2 & binary keystores, Actuator \texttt{/env} \\
    \bottomrule
  \end{tabular}
  \end{center}

  \vspace{0.6em}
  \begin{correction}{one item moved}
    Multi-line values were listed \unsolved. They work with one quoting rule,
    so the unsolved count is 2, not 3 --- and the ``keep PEM elsewhere''
    workaround is unnecessary.
  \end{correction}
\end{frame}

\begin{frame}{What is coming, in priority order}
  \small
  \begin{tabular}{@{}clp{5.6cm}@{}}
    \toprule
    & \textbf{Feature} & \textbf{Fixes} \\
    \midrule
    P0 & Scan an explicitly-named excluded path & \file{build/}, \file{dist/} \\
    1  & Monorepo workspaces --- \file{evnx.toml} listing app dirs, \cmd{validate -{}-all}, \cmd{cloud pull -{}-all} & three directories, three commands \\
    2  & Shared values / cross-vault references & one URL, three names \\
    3  & Escaped quotes in double-quoted multi-line values & the JSON parse failure \\
    4  & Prefix-aware validation (Vite, CRA) & \texttt{VITE\_STRIPE\_KEY} passing clean \\
    5  & \cmd{convert -{}-to spring-properties} / \texttt{application-yaml} & relaxed-binding naming \\
    6  & HashiCorp Vault / Spring Cloud Config targets & the overlap question \\
    \bottomrule
  \end{tabular}
\end{frame}

\begin{frame}[plain]
  \vfill
  \begin{center}
    {\LARGE\textbf{Questions}}\\[1.2em]
    {\small \url{https://www.evnx.dev/guides} · \url{https://github.com/urwithajit9/evnx}}\\[0.8em]
    {\footnotesize\textcolor{evnxMuted}{evnx 0.5.2 · every terminal block is a real run}}
  \end{center}
  \vfill
\end{frame}

\end{document}
